% context: 9th-grade physics homework, harder companion to 2-8PreTest_Kinematics for
%        early finishers; given 13 November 2026, before the unit test on 19 November
% topic: stopping distance (reaction plus braking, from km/h); a train catching a
%        cyclist (two position equations set equal); a stone thrown up from a roof
%        that lands below; the depth of a well from a dropped stone, then corrected
%        for the time sound takes to come back up
% structure: two pages, two problems per page, numbered 9-12 to follow the
%        PreTest's eight, with plenty of room to work
% design notes: no back-references to other sheets (values are restated); no
%        coincidental values within a problem (checked); multi-part problems build
%        step by step so each part is reachable
% data: braking a = -6.5 m/s^2 from the UK Highway Code 70 mi/h figure (2-9HW Q4);
%        reaction time 0.70 s is a chosen value; speed of
%        sound 343 m/s at 20 C (contexts.md #9); the train, roof and well are chosen
\documentclass[12pt, twoside]{article}
%\documentclass[12pt, twoside]{article}
\usepackage[letterpaper, margin=1in, headsep=0.2in]{geometry}
\setlength{\headheight}{0.6in}
%\usepackage[english]{babel}
\usepackage[utf8]{inputenc}
\usepackage{microtype}
\usepackage{amsmath}
\usepackage{amssymb}
%\usepackage{amsfonts}
\usepackage[nomessages]{fp} %\FPeval{\var-name}{2*sin(pi/6)}
\usepackage{siunitx} %units in math. eg 20\milli\meter
\usepackage{yhmath} % for arcs, overparenth command
\usepackage{tikz} %graphics
\usetikzlibrary{quotes, angles, arrows, arrows.meta}
\usepackage{graphicx} %consider setting \graphicspath{{images/}}
\usepackage{parskip} %no paragraph indent
\usepackage{enumitem}
\usepackage{multicol}
\usepackage{venndiagram}

\usepackage{fancyhdr}
\pagestyle{fancy}
\fancyhf{}
\renewcommand{\headrulewidth}{0pt} % disable the underline of the header
\raggedbottom
\hfuzz=2mm %suppresses overfull box warnings

\usepackage{hyperref}
\usepackage{float}

\title{Physics}
\author{Chris Huson}
\date{November 2026}

\fancyhead[LE]{\thepage}
\fancyhead[RO]{\thepage \\ Nickname: \hspace{2.25cm} \,\\ \,}
\fancyhead[LO]{La Scuola d'Italia / Huson / Physics \\* 13 November 2026}

\begin{document}

\subsubsection*{2.8 PreTest: Early Finishers (homework)}

\textbf{Reference (provided).} \quad $v = v_0 + at$ \quad
$x = x_0 + v_0 t + \frac{1}{2}at^2$ \quad $v^2 = v_0^2 + 2a(x - x_0)$

$\bar{v} = \frac{1}{2}(v + v_0)$ \quad $g = 9.80$ m/s$^2$ downward \quad 1 m/s = 3.6 km/h \quad Show your work, with units.

\vspace{0.2cm}

\begin{enumerate}[itemsep=0.4cm, start=9]

\item \textbf{Stopping distance.} A car is driving at 90 km/h when the driver sees
  a hazard ahead. It takes the driver 0.70 s to react and press the brake. During
  that time the car keeps its speed. Then the car brakes at $a = -6.5$ m/s$^2$.
  \begin{enumerate}[itemsep=0.4cm]
    \item How far does the car travel while the driver reacts?
      \vspace{2.0cm}
    \item Find the braking distance, and the total stopping distance.
      \vspace{3.4cm}
  \end{enumerate}

\item \textbf{Catching up.} A train pulls away from rest with $a = 0.40$ m/s$^2$.
  At the same moment a cyclist passes the front of the train, riding beside the
  track at a steady 6.0 m/s. Take the train's starting point as $x = 0$.
  \begin{enumerate}[itemsep=0.4cm]
    \item Write an equation for the position of each, the cyclist and the train,
      at time $t$.
      \vspace{1.3cm}
    \item When and where does the train catch up with the cyclist? \emph{Hint:
      their positions are equal.}
      \vspace{2.5cm}
    \item How fast is the train going at that moment? Compare it with the
      cyclist's speed, and explain why.
  \end{enumerate}

\newpage

\item \textbf{Off the roof.} A stone is thrown straight up at 12.0 m/s from the edge
  of a roof 30.0 m above the street. It misses the roof on the way down and lands
  in the street. Take \emph{up} as positive, with the roof edge at $y = 0$.
  \begin{enumerate}[itemsep=0.4cm]
    \item How high above the roof does the stone rise?
      \vspace{2.8cm}
    \item What is its velocity just before it hits the street? \emph{Hint: what is
      $y - y_0$?}
      \vspace{3.0cm}
    \item How long is the stone in the air?
      \vspace{2.4cm}
  \end{enumerate}

\item \textbf{How deep is the well?} You drop a stone down a well and hear the
  splash 2.40 s later.
  \begin{enumerate}[itemsep=0.4cm]
    \item Assume the sound reaches you instantly. How deep is the well?
      \vspace{2.4cm}
    \item As you recall, sound travels at 343 m/s in air at 20 \si{\celsius}. How
      long does the sound of the splash take to travel back up that depth?
      \vspace{1.8cm}
    \item So the stone fell for less than 2.40 s. Use the shorter time to estimate
      the depth again. Was it reasonable to ignore the time for the sound?
  \end{enumerate}

\end{enumerate}
\end{document}
